\section{Factoryn}
\label{sec:factory}

Alla delar finns nu. Det som fattas är limmet:

\begin{console}
EchoNode         behöver ett  Interface&
Spi              ÄR ett       Interface,  och behöver en ByteTransport&
AvrSpiTransport  ÄR en        ByteTransport
\end{console}

Någon måste skapa objekten, i rätt ordning, och hålla dem vid liv så länge programmet kör. Den
någon ska inte vara \code{main}, och definitivt inte \code{EchoNode}.

\begin{cppcode}
class Factory
{
public:
    Factory() noexcept;

    Factory(const Factory&)            = delete;
    Factory(Factory&&)                 = delete;
    Factory& operator=(const Factory&) = delete;
    Factory& operator=(Factory&&)      = delete;

    /**
     * @brief Get the CAN driver.
     *
     * @return Reference to the driver, valid for the lifetime of the factory.
     */
    [[nodiscard]] driver::can::Interface& canDriver() noexcept;

private:
    AvrSpiTransport myTransport; /**< Byte transport. Declared before the driver that uses it. */
    driver::can::Spi myDriver;   /**< The CAN driver, bound to myTransport. */
};
\end{cppcode}

\textbf{Deklarationsordningen är ett krav, inte stil.} Medlemmar konstrueras i den ordning de
deklareras, och \code{myDriver} tar en referens till \code{myTransport} i sin konstruktor. Byt
plats på de två raderna och du binder en referens till ett objekt som inte konstruerats än.
Kompilatorn säger ingenting.

\textbf{Returtypen är \code{Interface&}, inte \code{Spi&}.} Det är hela poängen. Anroparen får
något den kan använda, inte något den kan känna igen.

\textbf{Ingen heap.} Objekten är medlemmar, factoryn ligger på stacken i \code{main}, och allt har
statisk storlek. På ett friställt mål finns ingen \code{std::unique_ptr} att ta till --- och den
behövs inte heller. Smarta pekare i ett factory-exempel löser ett ägandeproblem som inte finns när
antalet objekt är känt vid kompileringstillfället.

\section{\texorpdfstring{\code{main}}{main}}
\label{sec:main}

\begin{cppcode}
int main()
{
    Factory factory{};
    app::EchoNode node{factory.canDriver(), ResponseId};

    while (true)
    {
        node.run();
    }
}
\end{cppcode}

Fem rader. \code{main} gör tre saker och inget mer: skapar factoryn, skapar applikationskomponenten
med det factoryn lämnar ut, och kör.

Orden \code{Spi}, \code{AvrSpiTransport}, \code{SPI0} och \code{PORTC} förekommer inte.
\textbf{Om du kan läsa din \code{main} utan att veta vilken hårdvara programmet kör på har du fått
arkitekturen rätt.}

\subsection{Testet på designen}

Byt en rad i factoryn:

\begin{cppcode}
private:
    driver::can::Stub myDriver;   // was: AvrSpiTransport + Spi
\end{cppcode}

Kompilera för \textbf{värddatorn}. Det ska gå, och programmet ska köra --- \code{EchoNode} ekar
frames mot en array i minnet i stället för mot en CAN-buss.

Går det inte, pekar kompilatorfelet ut exakt var någon rad ovanför interfacet
(\file{driver/can/interface.hpp}) har lärt sig något om hårdvaran som den inte borde veta.

Det är ett test du kan köra på tio sekunder, och det säger mer om arkitekturens kvalitet än någon
kodgranskning.

\section{En applikation som går att felsöka med}
\label{sec:debugapp}

\code{EchoNode} är bra att bevisa arkitekturen med och dålig att felsöka med: den gör ingenting
synligt. Inför \kapref{ch:bringup} är det värt att bygga en till.

\begin{narrowtable}{@{}ll@{}}
\thead{Lysdiod} & \thead{Tänds när} \\ \midrule
Grön & En frame har tagits emot och kvitterats. \\
Gul & En frame har skickats (\code{send()} returnerade \code{true}). \\
Röd & \code{hasError()} returnerade \code{true}. \\
\end{narrowtable}

Plus en knapp som skickar en frame med ett känt ID och en känd nyttolast.

Det är ett kvarts arbete och det betalar sig direkt vid bänken: en röd lysdiod som tänds när du
kopplar in transceivern säger något helt annat än en nod som inte gör något alls.

\textbf{Håll den utanför \code{EchoNode}.} Lysdioderna hör till applikationen, inte till
drivrutinen, och drivrutinen ska fortsätta vara byggbar för värddatorn.
